% Part: applied-modal-logics
% Chapter: epistemic-logic
% Section: introduction

\documentclass[../../../include/open-logic-section]{subfiles}

\begin{document}

\olfileid{aml}{el}{int}

\olsection{प्रस्तावना}

मोडल तर्कशास्त्र जसे \emph{मोडल विधानांचा} आणि
त्यांमधील तार्किक निष्पन्नतेच्या संबंधांचा विचार
करते, तसे ज्ञानविषयक तर्कशास्त्र \emph{ज्ञानविषयक
विधानांचा} आणि त्यांमधील निष्पन्नता संबंधांचा विचार
करते. मोडल संकारकांचा अर्थ संभाव्यता आणि
अनिवार्यता असा लावण्याऐवजी, एकस्थानी संयोजकांना
ज्ञानविषयक किंवा विश्वासविषयक अर्थ दिला जातो;
त्यायोगे ज्ञान आणि विश्वासाचे प्रतिमानीकरण करता
येते. उदाहरणार्थ, पुढील विधाने व्यक्त करायची
असू शकतात:

\begin{enumerate}
\item कॅलगरी हे अल्बर्टात आहे, हे रिचर्डला माहीत आहे.
\item सोफ्यावर कुत्रा असू शकतो, असे ऑड्रीला वाटते.
\item मंगळवारी तिचा वर्ग असतो, हे ऑड्रीला माहीत
  आहे, हे रिचर्डला माहीत आहे.
\item एका वर्षात १२ महिने असतात, हे सर्वांना माहीत आहे.
\end{enumerate}
आधुनिक ज्ञानविषयक तर्कशास्त्राचा उगम अनेकदा
याक्को हिंटिक्का यांच्या १९६२ मधील \emph{Knowledge
and Belief} या ग्रंथाशी जोडला जातो. तो ग्रंथ
लिहिला गेला तेव्हा संभाव्य जगांचे अर्थनिर्धारण
तर्कशास्त्रात वाढत्या प्रमाणात वापरले जात होते.
वास्तविक, ज्ञानविषयक तर्कशास्त्रे इतर मोडल
तर्कशास्त्रांसारखीच बहुतेक अर्थनिर्धारणाची साधने
वापरतात; मात्र त्यांचा अर्थ वेगळा लावतात. मुख्य
बदल म्हणजे \emph{प्राप्यता संबंध} कशाचे प्रतिनिधित्व
करतो, हे. ज्ञानविषयक तर्कशास्त्रांत तो काही प्रकारची
\emph{ज्ञानविषयक संभाव्यता} दर्शवतो. आपण कोणत्या
ज्ञानविषयक संकल्पनेचे प्रतिमानीकरण करतो त्यावर
प्राप्यता संबंधावर कोणत्या अटी घालायच्या हे अवलंबून
असेल, हे पुढे पाहू. अनुरूपता सिद्धांताला ज्ञानविषयक
अर्थ दिल्यावर काय घडते तेही पाहू. वरील उदाहरणांत
रिचर्ड आणि ऑड्री हे दोन कर्ते आहेत, आणि प्रत्येकीला
काय माहीत आहे यांतील संबंध येतो. आपण विचारात
घेणार असलेल्या ज्ञानविषयक तर्कशास्त्रांत अनेक कर्ते
असतील, म्हणून अशा गोष्टी व्यक्त करता येतील.
याउलट, एका कर्त्याचे ज्ञानविषयक तर्कशास्त्र फक्त
एकाच व्यक्तीला काय माहीत आहे किंवा काय वाटते
याबद्दल बोलेल.

\end{document}
